% Independently authored for FINMON_CP2.  Include in a document with amsthm.
\subsection{Finite monoids with bounded homology}

We count identities among the arrows of a finite category.  In this
subsection, a finite monoid is regarded as a category with one object.
The universal group of a monoid $M$ is denoted by $U(M)$.
The checks below confirm the four assertions requested in C.1.  The
chronology of the independent derivations is recorded separately in the
check-event table; a later comparison with the supplied proofs is not
mislabelled as an independent derivation.

\begin{lemma}[Proved: nilpotence from a cohomological bound]
\label{CP2:C-Knilpotence}
Let $R=\mathbb Z$ or $R=\mathbb Z_{(p)}$, and let $X$ be a connected CW
complex such that $H^q(X;R)=0$ for all $q>n$.  The ideal of classes of
rank zero in $\mathrm{KU}_R^0(X)$ is nilpotent.
\end{lemma}
\begin{proof}
We give the convergence argument explicitly.  Let $Y$ be the zero
component of the zeroth infinite-loop space of $\mathrm{KU}_R$.
Its positive homotopy groups are $R$ in even degrees and zero in odd
degrees.  Choose an even integer $N>n+1$, and let $F$ be the homotopy
fiber of the Postnikov map $Y\to P_NY$.
The finite Postnikov sections of $F$ have their nonzero homotopy groups
in even degrees greater than $N$.  Obstruction theory therefore gives
\[
 [X,P_mF]=0,\qquad [\Sigma X,P_mF]=0
\]
for every $m$: the relevant obstruction and choice groups are ordinary
cohomology groups in degrees greater than $n$.  We use reduced
cohomology and based maps here.  Since a simply connected space is the
homotopy limit of its Postnikov sections, the mapping-space homotopy
limit sequence has zero inverse-limit and zero derived-limit terms in
these two degrees.  It follows that $[X,F]=0$ and hence that
\[
 [X,Y]\longrightarrow [X,P_NY]
\]
is injective.  This argument applies to arbitrary CW complexes $X$;
finite type is not assumed.

The product of $r$ rank-zero classes is induced by the reduced
multiplication map $Y^{\wedge r}\to Y$.  The factorization through the
smash product follows because a zero factor makes the product zero;
it can equally be expressed using the multiplication of the ring
spectrum $\mathrm{KU}_R$.  As $Y$ is simply connected,
$Y^{\wedge r}$ is $(2r-1)$-connected.  If $2r-1\ge N$, its map to the
$N$-type $P_NY$ is nullhomotopic.  The product is consequently zero by
the injectivity just proved.  The same $r$ works for every choice of
the rank-zero classes.
\end{proof}

\begin{theorem}[Proved: R9 Theorems 1.1 and 1.2]
\label{CP2:C-finite-group-maps}
Let $X$ be a connected CW complex with finite fundamental group
$\pi$, let $G$ be a finite group, and let $f:BG\to X$ be a continuous
map.  If $H^q(X;\mathbb Z)=0$ in all sufficiently high degrees, then
$f_*:G\to\pi$ is trivial.  More precisely, if
$H^q(X;\mathbb Z_{(p)})=0$ in all sufficiently high degrees, then
$p$ does not divide the order of $f_*(G)$.
\end{theorem}
\begin{proof}
Let $c:X\to B\pi$ classify the universal covering.  The regular
complex representation of $\pi$ gives a vector bundle $V\to X$ of
rank $d=|\pi|$.  By Lemma~\ref{CP2:C-Knilpotence},
$v=[V]-d$ is nilpotent in the corresponding K-theory ring.

Suppose first that integral cohomology is bounded and $f_*(G)$ is
nontrivial.  There is an element $g\in G$ of prime-power order
$p^a$ with $f_*(g)\ne1$: take a preimage of an element of order $p$
and then take its $p$-primary component.  Put $P=\langle g\rangle$.
The composite $BP\to BG\xrightarrow{f}X$ pulls $v$ back to the
virtual representation
\[
 w=\operatorname{Res}_{P}(\mathbb C[\pi])-d.
\]
Atiyah's Theorem~7.2 identifies the inverse limit of K-theory over
finite skeleta of $BP$ with the augmentation-ideal completion of
$R(P)$; Proposition~6.11 says that $R(P)$ injects into that
completion~\cite{CP2-Atiyah1961}.  A zero power of the pulled-back
class would therefore give a zero power of $w$ in $R(P)$.  Evaluation
of its character at $g$, however, gives $w(g)=-d$, so no power is
zero.  This is a contradiction.

For the prime-local assertion, suppose $p$ divides $|f_*(G)|$ and
choose $P$ as above for this prime.  We avoid commuting localization
with an infinite inverse limit.  Model $BP$ by the infinite lens
space and restrict to each finite odd-dimensional lens skeleton
$L^{2j+1}$.  Its reduced integral $K^0$ is a finite $p$-group: the
finite Atiyah--Hirzebruch spectral sequence has only finite
$p$-groups in the positive even positions contributing to $K^0$.
Thus
\[
 K^0(L^{2j+1})\longrightarrow
 K^0(L^{2j+1})\otimes\mathbb Z_{(p)}
\]
is injective.  On a finite CW complex the latter is
$\mathrm{KU}_{(p)}^0$.  Nilpotence of the localized pullback therefore
makes one fixed power of $w$ vanish integrally on every such skeleton.
Atiyah's Theorem~7.2 and Proposition~6.11 again force that power to
vanish in $R(P)$, contradicted by its character value $(-d)^r$.

Every map and bundle pullback in this proof is in the category of
based CW spaces up to based homotopy.  Removing basepoints only
conjugates the induced group homomorphism and does not change either
conclusion.
\end{proof}

\begin{remark}
\label{CP2:C-linear-extension}
The regular representation can be replaced by any finite-dimensional
complex representation of $\pi_1(X)$.  A finite-order complex matrix
is diagonalizable with eigenvalues on the unit circle, so its trace
equals its dimension only when it is the identity.  Consequently,
without assuming $\pi_1(X)$ finite, the image of $f_*$ lies in the
intersection of the kernels of all such representations.  In
particular the integral assertion holds when $\pi_1(X)$ is
residually finite or is a complex linear group.
\end{remark}

\begin{lemma}[Proved: the finite-group corner and its comparison map]
\label{CP2:C-minimal-corner}
For a finite monoid $M$ there is an idempotent $e$ such that $H=eMe$
is a finite group and a map $j_e:BH\to BM$ inducing a surjection
$H\to U(M)$ on fundamental groups.  In particular $U(M)$ is finite.
\end{lemma}
\begin{proof}
Choose a minimal nonempty two-sided ideal $I$ of $M$ and an idempotent
$e\in I$; a suitable positive power of any element of $I$ is
idempotent because its powers are eventually periodic.  For
$x\in eMe$, minimality gives $MxM=I$, hence
$e=axb$ for some $a,b\in M$.  Thus
\[
 e=(eae)x(ebe)
\]
inside the finite monoid $eMe$.  In a finite monoid, a product equal
to the identity forces both factors to be units: if $uv=1$, left
multiplication by $u$ is surjective and hence injective, so
$vu=1$.  Applying this twice to the displayed factorization shows
that every $x\in eMe$ is a unit.  Thus $H=eMe$ is a group.

The universal map $\eta:M\to U(M)$ sends idempotents to $1$, and
\[
 \eta(m)=\eta(eme) \qquad(m\in M).
\]
The subgroup $\eta(H)$ therefore contains every generator of
$U(M)$, proving surjectivity.

We spell out the topological comparison, since the inclusion
$eMe\subset M$ need not preserve the identity.  Let
$\operatorname{Kar}(M)$ have the idempotents as objects and
$\operatorname{Kar}(M)(a,b)=bMa$; composition is multiplication.
The full inclusion $i:M\to\operatorname{Kar}(M)$ at the object $1$
is a homotopy equivalence on classifying spaces.  Indeed, for an
idempotent $a$, the comma category $(i\downarrow a)$ has objects
$x\in aM$ and morphisms $m:x\to y$ satisfying $ym=x$.
The arrows $x:x\to a$ define a natural transformation from the
identity functor to the constant functor at $a$; naturality is
exactly $ym=x$.  This contracts its classifying space.  Quillen's
Theorem~A proves that $Bi$ is a homotopy equivalence.
The full inclusion of the object $e$ now defines the span
\[
 BH\longrightarrow B\operatorname{Kar}(M)
       \xleftarrow[\simeq]{Bi} BM.
\]
We define $j_e$ using a homotopy inverse of $Bi$.  Its effect on
fundamental groups is the homomorphism $\eta|_H$, since the two
arrows represented by $e$ identify the object $e$ with the object
$1$ after groupoid completion.  This proves the asserted typing and
surjectivity of the map.
\end{proof}

\begin{theorem}[Proved: R9 Theorem 2.1]
\label{CP2:C-unbounded-homology}
If a finite monoid $M$ has $H_k(BM;\mathbb Z)=0$ in all sufficiently
high degrees, then $U(M)=1$.  If a prime $p$ divides $|U(M)|$, then
$H_k(BM;\mathbb F_p)$ is nonzero in infinitely many degrees.
\end{theorem}
\begin{proof}
The integral universal-coefficient exact sequence converts eventual
vanishing of integral homology into eventual vanishing of integral
cohomology.  Apply Theorem~\ref{CP2:C-finite-group-maps} to the map
of Lemma~\ref{CP2:C-minimal-corner}; its induced homomorphism is both
surjective and trivial, so $U(M)=1$.

For the second assertion, suppose mod-$p$ homology vanished in all
large degrees.  The normalized bar chain group in every degree has
finite rank because $M$ is finite, so all integral homology groups
are finitely generated.  The homology universal-coefficient exact
sequence implies $H_k(BM;\mathbb Z)\otimes\mathbb F_p=0$ for all
large $k$.  A finitely generated abelian group with this property
has zero localization at $p$.  The cohomology universal-coefficient
sequence then gives eventual vanishing of
$H^q(BM;\mathbb Z_{(p)})$.  The prime-local conclusion of
Theorem~\ref{CP2:C-finite-group-maps}, applied to the same surjection,
implies $p\nmid|U(M)|$, a contradiction.
\end{proof}

\begin{theorem}[Proved: R9 Theorem 2.2]
\label{CP2:C-finite-equivalences}
For a finite monoid $M$, the following are equivalent:
\begin{enumerate}
\item $BM$ is homotopy equivalent to a finite CW complex;
\item $BM$ is homotopy equivalent to a finite-dimensional CW complex;
\item $H_k(BM;\mathbb Z)=0$ in all sufficiently high degrees.
\end{enumerate}
When these conditions hold, $BM$ is simply connected.  If homology
vanishes above degree $d$, a finite model can be chosen of dimension
at most $d+1$; if $H_d$ is torsion-free, dimension $d$ suffices.
\end{theorem}
\begin{proof}
The first implication and the implication from the second condition
to the third follow from cellular homology.  Under the third
condition Theorem~\ref{CP2:C-unbounded-homology} gives simple
connectivity.  The bar complex has finite rank in each degree, so
integral homology is finitely generated.  Hatcher's
Proposition~4C.1~\cite{CP2-Hatcher2002}, whose exact statement we
opened, supplies a homotopy-equivalent CW complex with one
generator cell for each infinite cyclic homology summand and a
generator--relator pair in successive dimensions for each finite
cyclic summand.  Only finitely many cells are required here.  The
dimension assertions follow from that same cell count.

An independent derivation recorded before reading the supplied
proof also constructs a finite model directly.  For $d\ge2$,
take $L=(BM)^{(d+1)}$.  Relative Hurewicz for $(BM,L)$ identifies
$H_{d+1}(L)$ with $\pi_{d+2}(BM,L)$.  The former is finite free.
Attach one $(d+2)$-cell for each basis element, using representative
relative disks in $BM$.  The resulting finite simply connected
complex maps to $BM$ by an integral homology isomorphism, and the
homological Whitehead theorem makes it a homotopy equivalence.
The Hatcher argument above gives the sharper stated dimension.
\end{proof}

\begin{corollary}[Proved using the stated outside theorem: R9 Corollary 2.5]
\label{CP2:C-characterization}
The homotopy types of classifying spaces of finite monoids having
bounded integral homology are exactly the homotopy types of finite
simply connected CW complexes.  The realization in the reverse
direction can be chosen to be aperiodic.
\end{corollary}
\begin{proof}
The forward implication is Theorem~\ref{CP2:C-finite-equivalences}.
For the reverse implication we use Ben-Bassat's Theorem~C
\cite{CP2-BenBassat2026}.  It realizes every finite simply connected
CW complex by a finite aperiodic monoid.  The resulting homology is
bounded because homotopy equivalence preserves homology.

We checked the proof at the steps used here.  Theorem~A identifies
the framed collage with the homotopy cofiber, respecting its map
from the old corner.  Lemma~5.3 replaces a finite connected
simplicial source by a reduced source through contractible
edge-replacement diagrams.  Theorem~B therefore applies to each
attaching sphere.  Theorem~C starts with a finite simply connected
CW model without one-cells and applies Theorem~B successively; each
attaching sphere is connected.  The aperiodicity check is the
stationary-power argument in Lemma~5.2.  We found no fault in these
steps.  We cite the unrefereed theorem as an outside result and do
not claim an independent re-proof of that whole paper.
\end{proof}

\begin{proposition}[Refuted: R8 Conjecture 7.1; proved replacement]
\label{CP2:C-conjecture-counterexample}
Finite fundamental group does not suffice for a connected finite CW
complex to be a finite-monoid classifying space.  The exact
replacement, using the outside theorem in
Corollary~\ref{CP2:C-characterization}, is simple connectivity.
In particular, $\mathbb{RP}^2$ is not homotopy equivalent to $BM$
for any finite monoid $M$.
\end{proposition}
\begin{proof}
The usual CW structure on $\mathbb{RP}^2$ has one cell in each of
dimensions $0,1,2$, with the two-cell attached along the square of
the one-cell loop.  It is finite and connected, its fundamental
group is $\langle a\mid a^2=1\rangle\cong C_2$, and its homology
vanishes above degree two.  Theorem~\ref{CP2:C-unbounded-homology}
therefore rules out every finite monoid.

For a completely enumerated finite counterexample we also give the
simplicial complex $K$ on vertices $1,\ldots,6$ with facets
\[
\begin{gathered}
123,\quad124,\quad135,\quad146,\quad156,\\
236,\quad245,\quad256,\quad345,\quad346.
\end{gathered}
\]
Its complete graph is connected.  Collapse the tree with edges
$12,13,14,15,16$.  Its fundamental-group presentation has the ten
remaining edges as generators and the ten triangular boundary
words as relators.  Nine elementary Tietze eliminations, recorded
literally in \texttt{rp2\_certificate.json}, leave one generator
with its square as the sole relator.  Thus $\pi_1(K)=C_2$.
The same certificate contains its cellular boundary matrix after
tree collapse; its determinant is $-2$, giving
$H_1(K;\mathbb Z)=\mathbb Z/2$ and $H_2(K;\mathbb Z)=0$.
All higher homology vanishes by dimension.  Our independent checker
verifies each Tietze substitution, the determinant by integer
elimination, the two-element group table, and the full composition
table of the face-poset category.  That category has 31 objects and
121 arrows including identities; its nerve is the barycentric
subdivision of $K$.  In particular this is a finite explicit
counterexample without any reliance on recognizing a surface.
Necessity and sufficiency of the corrected simple-connectivity
condition are precisely
Corollary~\ref{CP2:C-characterization}.
\end{proof}
